\section{The exact central Northcott reduction}
\label{sec:central-northcott}

Continue with the central notation of
\MathSecref{sec:central-projector}. Let
\(\Lambda_U(E)=\Pi_*\Lambda_Y^c(E)\) and let
\(\Lambda_G(E)\) be its image in \(G_E\). The quotient may be
viewed in the geometric group: if an element defined over \(E\)
becomes geometrically algebraically trivial, it is in \(A_E\) by
definition. In particular, these rational quotients inject under
field extension.

\begin{theorem}\label[theorem]{thm:central-northcott}
The following conditions are equivalent:
\begin{enumerate}[label=(\roman*)]
\item The form \(q_U(z)=(-1)^ch_Y(z,z)\) has bounded-degree
Northcott on \(\Lambda_U\).
\item For every finite \(E/F_0\), \(K_E=\ker(\AJ_\Q|_{A_E})=0\),
and \(q_G\) has bounded-degree Northcott on \(\Lambda_G\).
\end{enumerate}
There is also a fixed-field version, with all groups and the
vanishing condition taken over one fixed \(E\).
\end{theorem}

\begin{proof}
Suppose (i). If \(0\ne z\in K_E\), multiply by one integer to
make \(z\) integral. All positive integer multiples of this class
are distinct and have height zero, contradicting Northcott.
Thus \(K_E=0\).

Choose an integer \(M\ge1\) such that \(MR\) is represented by
an integral correspondence. For \(\xi\in\Lambda_G(E)\), choose
\(z\in\Lambda_U(E)\) mapping to \(\xi\). Then
\[
 MRz\in\Lambda_U(E),\qquad
 q_U(MRz)=M^2q_G(\xi),\qquad [MRz]=M\xi.
\]
Northcott for \(\Lambda_U\) bounds the possible classes \(MRz\),
and hence the possible \(\xi\). This proves residual Northcott.

Conversely, assume (ii). Residual Northcott implies
\(q_G(\xi)>0\) for every nonzero rational class \(\xi\in G_E\).
Indeed one multiple of \(\xi\) is in \(\Lambda_G(E)\), and a
nonpositive height would give infinitely many integral multiples
of height at most zero. For every \(z\in U_E\), orthogonality gives
\begin{equation}\label{eq:central-norm-split}
 q_U(z)=q_e(ez)+q_G([z]).
\end{equation}
Both terms are nonnegative. If \(z\) is integral and
\(q_U(z)\le B\), residual Northcott bounds the possibilities for
\([z]\). The denominators of the fixed projector \(e\), followed
by \Cref{lem:visible-integral}, bound the possibilities for \(ez\).
Finally
\[
 z\longmapsto(ez,[z])
\]
is injective: its kernel is \(RA_E=K_E=0\). This proves (i).
The argument with \(E\) fixed is identical.
\end{proof}

The equivalence separates the algebraic Abel--Jacobi kernel from
integral finiteness on the residual quotient. In fact, Northcott
also forces real positive definiteness, by
\Cref{lem:northcott-real}. This observation will combine the
central equivalence with the primitive calculation in
\Cref{thm:program-equivalence}.
